\documentclass[11pt,letterpaper]{article}
\usepackage[margin=0.94in]{geometry}
\usepackage[T1]{fontenc}
\usepackage{lmodern}
\usepackage{amsmath,amssymb,amsthm,mathtools}
\usepackage{booktabs,longtable,array}
\usepackage[dvipsnames]{xcolor}
\usepackage{microtype}
\usepackage{enumitem}
\usepackage{fancyhdr}
\usepackage{hyperref}
\usepackage{xurl}
\hypersetup{colorlinks=true,linkcolor=MidnightBlue,urlcolor=MidnightBlue,citecolor=MidnightBlue,pdftitle={Common remainder certificates and explicit Diophantine forms of Con(ZFC)},pdfauthor={OpenAI ChatGPT, prepared for ProveIt}}
\pagestyle{fancy}\fancyhf{}\fancyhead[L]{\small Common remainder certificates}\fancyhead[R]{\small ProveIt research note}\fancyfoot[C]{\thepage}
\setlength{\headheight}{14pt}
\setlength{\parskip}{0.35em}
\setlength{\parindent}{0pt}
\setlist{itemsep=2pt,topsep=5pt}
\newtheorem{theorem}{Theorem}[section]
\newtheorem{lemma}[theorem]{Lemma}
\newtheorem{proposition}[theorem]{Proposition}
\newtheorem{corollary}[theorem]{Corollary}
\theoremstyle{remark}\newtheorem{remark}[theorem]{Remark}
\newcommand{\Z}{\mathbb Z}\newcommand{\N}{\mathbb N}
\newcommand{\Con}{\operatorname{Con}}\newcommand{\ZFC}{\mathsf{ZFC}}
\newcommand{\NN}{\mathbb Z_{>0}}
\newcommand{\pc}{\operatorname{pc}}
\newcommand{\htp}{\operatorname{ht}}
\newcommand{\file}[1]{\nolinkurl{#1}}
\newcommand{\sourcepin}{\nolinkurl{55d1f72aedebf892e446c90d35fa1ae37689b0f2}}
\allowdisplaybreaks
\begin{document}
\hypersetup{pageanchor=false}
\begin{titlepage}
\vspace*{0.6in}
{\Large\color{MidnightBlue} DIOPHANTINE CONSISTENCY}\par
\vspace{0.35in}
{\huge\bfseries Common remainder certificates\par}
\vspace{0.13in}
{\LARGE and explicit polynomial forms of $\Con(\ZFC)$\par}
\vspace{0.4in}
{\large A quartic in 252 unrestricted integer variables,\par
smaller native coefficients, and improved Pell carriers}\par
\vspace{0.45in}
\begin{tabular}{@{}ll@{}}
Prepared for & Vladimir Reshetnikov / ProveIt\\
Date & 25 October 2026\\
Preparation & OpenAI ChatGPT; AI-assisted research\\
Source snapshot & \texttt{55d1f72aedeb}\\
\end{tabular}\par
\vspace{0.45in}
\begin{minipage}{0.94\linewidth}
\textbf{Main result.} A fixed, completely materialized polynomial
\[
 P_{252}(z_0,\ldots,z_{251})=\sum_{j=0}^{181}Q_j(z_0,\ldots,z_{251})^2
\]
has degree four and 5,340 collected monomials. Every $Q_j$ is explicitly
quadratic or affine. Composing the new arithmetic proof with the pinned
direct-calculus calibration establishes
\[
 \ZFC\vdash
 \bigl[\neg\exists\mathbf z\in\Z^{252}\;P_{252}(\mathbf z)=0\bigr]
 \longleftrightarrow \Con(\ZFC).
\]
The new step replaces eleven separate interval certificates by one common
quadratic bound on rescaled remainder coordinates. Relative to the newest
matched direct parent, it removes 30 variables and ten residual equations,
without increasing the degree or coefficient height.
\end{minipage}
\vfill
\textbf{Scope.} This is an ordinary mathematical proof with exact finite
algebraic verification. It does not assert a global smallest equation,
a proof of consistency, or a new end-to-end proof-assistant derivation.
The inherited logical and arithmetic dependencies are identified and
retained in the accompanying archive.
\end{titlepage}
\hypersetup{pageanchor=true}

\begin{abstract}
We construct a fixed integer quartic whose unsolvability is provably
equivalent in ZFC to the syntactic consistency of ZFC. The construction
uses the repository's direct numerical proof calculus, and does not encode
K215 or specialize Jones's universal polynomial. Its 252 unrestricted
integer variables improve the latest matched direct construction's 282.
The mechanism is a general representation lemma: digit remainders with
different positive divisors are rescaled to one common divisor, and their
separate nonnegativity certificates are replaced by one small Euclidean
norm certificate. A two-stage coefficient bound recovers the factorial
digit before its factorial interpretation is available, and then recovers
the ten range digits. We give constructive soundness and completeness,
the full residual specification, the complete collected polynomial, and
independent finite checks. Separate results complete a previously
conditional radix-$2^{49}$ native coefficient construction and lower the
degrees of compact ten-positive-variable equations by reusing a smaller
carrier denominator. A complemented-label encoding also permits a fully certified radix-$2^{47}$
native tuple, with 207,214 raw coefficient bytes at the unchanged
15-variable, degree-157, 66-operation profile.
\end{abstract}
\tableofcontents
\newpage

\section{What is constructed and what is being optimized}

The request is for an ordinary polynomial with fixed integer coefficients,
not an expression that leaves exponentiation, a proof predicate, or a
consistency-dependent constant unevaluated. The principal deliverable is
\file{direct/equation/equation.json}, whose variable order and 182 residuals define
$P_{252}$ exactly. The companion
\file{direct/equation/expanded_polynomial.json} lists every collected coefficient.
Appendix~\ref{app:residuals} prints the same residual specification using
the fixed variables $z_0,\ldots,z_{251}$.

The relevant comparison is multidimensional. The number of variables,
degree, coefficient height, expanded support, arithmetic-circuit cost,
and length of the coefficient specification are different measures.
Reducing one can increase another. The repository already contains
equations with fewer variables than 252 on other routes and at other
degrees; this report does not erase those alternatives. Its principal
improvement is within the direct proof-calculus quartic route, which has
small literal coefficients and needs neither a native machine tuple nor
a universal-polynomial specialization.

All comparisons use the fixed repository commit \sourcepin. The collection
index, active constructions, and newest matched incoming parents were
inspected. In particular, the direct matched parent is the
\file{strict_dimension_centered} construction in the incoming
\emph{Cofinal Anchors and Height One, Round 17} archive~\cite{round17}.
Comparing only against the older collection index would miss this parent.

\begin{table}[ht]
\centering
\caption{Matched direct quartics over unrestricted integers. Coefficient
height is the maximum absolute collected coefficient.}
\begin{tabular}{@{}lrr@{}}
\toprule
Measure & Pinned parent & New $P_{252}$\\
\midrule
Variables & 282 & \textbf{252}\\
Quadratic or affine residuals & 192 & \textbf{182}\\
Total degree & 4 & 4\\
Nonzero collected monomials & 5,491 & \textbf{5,340}\\
Coefficient height & 128,446 & 128,446\\
\bottomrule
\end{tabular}
\end{table}

Thus the new construction saves $30/282\approx10.64\%$ of the variables,
$10/192\approx5.21\%$ of the residuals, and $151/5491\approx2.75\%$ of
the expanded monomials. These are finite counts of the supplied objects,
not estimates from a symbolic template.

The general DPRM theorem explains why a Diophantine representation of a
recursively enumerable refutation predicate exists~\cite{matiyasevich}.
It does not supply these particular small explicit coefficients. Likewise,
Jones's variable reduction and later universal-pair work provide important
comparison frameworks~\cite{jones,bayer}; they are not substituted for
the direct construction or its calibration here.

\section{A common certificate for many remainders}\label{sec:lemma}

For integers $a,b,c$, define
\[
 G(a,b,c)=a^2+b^2+c(c+1).
\]
This expression is always nonnegative. Conversely, every $n\in\N$ has
such a representation. Indeed, $4n+1$ is not of the form
$4^u(8v+7)$, so the three-square theorem represents it by three squares.
Reduction modulo four shows that exactly one of those squares is odd.
Writing the odd root as $2c+1$ and the other roots as $2a,2b$ gives
$n=G(a,b,c)$. This is the same elementary integer-sign representation
used by the parent.

\begin{lemma}[Integral rescaling of a remainder]\label{lem:rescale}
Let $L,M$ be positive integers with $L\mid M$, and put $v=M/L$.
For integer $U,Q,R$, the equation
\[
 vU=MQ+R
\]
implies $v\mid R$. The integer $r=R/v$ then satisfies $U=LQ+r$.
Conversely every integer solution of the latter equation gives a solution
of the former by $R=vr$.
\end{lemma}
\begin{proof}
Both $vU$ and $MQ=LvQ$ are divisible by $v$, so their difference $R$ is
divisible by $v$. Division by the nonzero integer $v$ gives the assertion.
The converse is direct multiplication. In the polynomial system $v$ is
an integer coordinate constrained by $Lv=M$; no rational coordinate is
introduced.
\end{proof}

\begin{theorem}[One norm certificate for a family of small remainders]
\label{thm:common}
Let $M,L_1,\ldots,L_s$ be positive integers with $L_i\mid M$, and
$v_i=M/L_i$. Suppose the scaled digit equations
\[
 v_iU_i=MQ_i+R_i \qquad(1\le i\le s)
\]
hold, together with
\begin{equation}\label{eq:norm}
 M^2=1+4\sum_{i=1}^sR_i^2+G(a,b,c).
\end{equation}
Then $r_i=R_i/v_i$ is an integer with $|r_i|<L_i/2$. If the ordinary
least nonnegative remainder $r_i^0$ of $U_i$ modulo $L_i$ satisfies
$r_i^0<L_i/2$, then $r_i=r_i^0$.

Conversely, if $s\le11$ and an ordinary digit family has
$0\le r_i^0<L_i/32$, there exist $a,b,c\in\Z$ for which its scaled
remainders $R_i=v_i r_i^0$ satisfy~\eqref{eq:norm}.
\end{theorem}
\begin{proof}
Lemma~\ref{lem:rescale} gives integrality. Nonnegativity of $G$ implies
$4R_i^2<M^2$, hence $|r_i|<L_i/2$. The difference $r_i-r_i^0$ is a
multiple of $L_i$ and has absolute value strictly less than $L_i$; it is
therefore zero.

For the converse, $0\le R_i<M/32$, and so
\[
 4\sum_iR_i^2<\frac{44}{1024}M^2
 =\frac{11}{256}M^2<M^2.
\]
The difference $M^2-4\sum_iR_i^2$ is a positive integer, so subtracting
one leaves a nonnegative integer. Represent it by $G$.
\end{proof}

The completeness hypothesis is essential: a common norm need not accept
every arbitrary tuple of remainders in their individual intervals.
Here it accepts the much smaller canonical remainders produced by the
binomial constructions. The next two sections prove that margin without
assuming the conclusion of the factorial step.

\begin{remark}[A valid smaller coefficient]
The coefficient four in~\eqref{eq:norm} can be replaced by two when the
$1/32$ canonical bound is available in both directions. The new bound is
$|r_i|<L_i/\sqrt2<3L_i/4$. Thus
$|r_i-r_i^0|<25L_i/32<L_i$, which still forces equality. Completeness
uses $22M^2/1024<M^2$. The released primary polynomial uses four for the
particularly transparent centered interval. No second numerical endpoint
is claimed from this observation.
\end{remark}

\section{Applying the lemma to the direct arithmetic source}

The inherited source uses eighteen certified power requests. These are
implemented by ordinary polynomial Pell relations. A symbol such as
$q^t$ below denotes the integer value of an existing output coordinate
after its meaning has been proved; it is not an operation left inside
the final polynomial.

The eleven replaced remainder intervals have only four different
denominators:
\begin{center}
\begin{tabular}{@{}lll@{}}
\toprule
Denominator & Number & Purpose\\
\midrule
$q^t$ & 2 & Factorial binomial; encoded index\\
$q^{t+1}$ & 5 & Four syntax representatives; derived label\\
$q^{y+1}$ & 2 & Two pending-multiset representatives\\
$q^x$ & 2 & Two beta quotients\\
\bottomrule
\end{tabular}
\end{center}

The parent proves
\[
 t\ge16,\quad y\ge1,\quad x\ge2+(t+1)y
\]
before any of these eleven interval guards is used. Set $M=q^x$ and
introduce three ratio coordinates by
\begin{equation}\label{eq:ratios}
 q^t v_t=M,\qquad q^{t+1}v_{t+1}=M,\qquad q^{y+1}v_y=M.
\end{equation}
All the divisors in these equations are positive powers of the same
positive base, with exponents at most $x$, so the ratios are positive
integers. For the fourth denominator the ratio is simply one.

The factorial digit formerly had the form
\[
 U=q^t(B+q\omega)+r.
\]
It becomes
\[
 v_tU=M(B+q\omega)+R.
\]
The quantity $q\omega$ already has a circuit coordinate, so this is a
quadratic residual in the actual coordinates. Each range digit formerly
had the form
\[
 U=Nq^k\omega+r
\]
and becomes $v_kU=NM\omega+R$. The four separate product coordinates
$Nq^k$ are consequently replaced by one shared coordinate $NM$.
The new remainder coordinates occur only in their scaled digit equations
and the one norm equation; they do not change a private Pell frame atom
or a requested power.

\subsection{The order in which signs and power meanings are recovered}

The new proof retains the parent's signed-frame and semantic-anchor
arguments, followed by its coordinate-orientation and strict-budget
arguments. Their conclusions are used in the following order:

\begin{enumerate}
\item The three retained initial sign predicates imply $B\ge1$ and
      $q\ge B+1\ge2$. The cofinal horizon has $t\ge16$.
\item Signed Pell orientation and the request at exponent $t$ identify
      $V_0=(F^2+1)^t$. The raw ceiling $FB=V_0+s$ with $s\ge0$ then
      gives $F\ge1$ and $B>F$, hence $B\ge2$.
\item The beta value $y=V_0(V_0-1)$ and the two endpoint equations give
      $x\ge2+(t+1)y\ge19$.
\item The requests at base $q+1$ exclude negative actual representative
      indices. Thus the nine actual CRT representatives are nonnegative.
      Writing $H=1+\sum a_i$ gives $H\ge1$.
\item The remaining requests identify $N=B^x$ and
      $q=(BH)^x=B^xH^x$. Every denominator and numerator power is now
      interpreted.
\end{enumerate}

Only after these steps do we interpret~\eqref{eq:ratios} and the common
norm. In particular, no deleted remainder sign is needed to establish
the positive divisor or power used to recover that sign. The inherited
proofs are included under \file{direct/upstream/round17/}; their boundary
is discussed in Section~\ref{sec:calibration}.

\section{The factorial digit is recovered before assuming $F=t!$}

\begin{lemma}[An elementary coefficient bound]\label{lem:binom}
For integers $m,t\ge1$ and $0\le j\le t$,
\[
 \binom{m+t}{j}\le(m+1)^t.
\]
\end{lemma}
\begin{proof}
For $t=1$ the assertion is immediate. In the next Pascal row, the
coefficients with $1\le j\le t$ are bounded by
$2(m+1)^t\le(m+1)^{t+1}$; the constant edge is one. The newly included
last coefficient satisfies
\[
 \binom{m+t+1}{t+1}
 =\prod_{i=1}^{t+1}\left(1+\frac mi\right)
 \le(m+1)^{t+1}.
\]
This completes the induction.
\end{proof}

Apply the lemma with $m=F^2$. The already interpreted factorial numerator
is $U=(q+1)^{F^2+t}$, and every binomial coefficient through position $t$
is at most $V_0=(F^2+1)^t$. The raw ceiling, before factorial semantics,
gives
\[
 V_0\le BF<B^2.
\]
Because $q=B^xH^x$ and $x\ge8$,
\begin{equation}\label{eq:rawbound}
 \frac{V_0}{q-1}
 <\frac{B^2}{q-1}
 <\frac{2}{B^{x-2}}\le\frac1{32}.
\end{equation}

The true low part in the binomial expansion is
\[
 r^0=\sum_{j=0}^{t-1}\binom{F^2+t}{j}q^j.
\]
The geometric-series bound and~\eqref{eq:rawbound} give
\[
 0\le r^0<\frac{V_0q^t}{q-1}<\frac{q^t}{32}.
\]
In particular this is indeed the ordinary remainder modulo $q^t$.
Theorem~\ref{thm:common} reconstructs it exactly from the new signed
remainder. Dividing the restored digit equation by $q^t$ and reducing
modulo $q$ gives
\[
 B\equiv\binom{F^2+t}{t}\pmod q.
\]
Both quantities lie in $[0,q)$: $B<q$ follows from the retained initial
output slack, and the binomial coefficient is at most $V_0<q$.
Consequently
\[
 B=\binom{F^2+t}{t}.
\]
The unchanged factorial fixed-point theorem of the direct source now
gives $F=t!$. The proof above has not used this identity early.

\section{The remaining range digits and constructive completeness}

After the factorial interpretation, the inherited arithmetic congruence
$B\mid C_0+F^2+1$ excludes a zero encoded-index representative: the
factorial formula gives $B>F^2+1$. Thus $C_0\ge1$ and $H\ge2$.

For an ordinary nonnegative upper representative $a$ and a relevant
position $0\le j\le k\le x$,
\[
 \binom aj\le\max(1,a)^k\le H^x.
\]
This includes $a=j=0$; using $a^k$ without the maximum would give a
false bound at that edge. For the derived label $h$, the inherited
literal degree bound is $h\le H^\ell$ with $\ell=4,194,304$, and the
factorial/endpoint estimate gives $\ell(t+1)<x$. Its relevant
coefficients therefore have the same $H^x$ upper bound.

Since $q=B^xH^x$,
\[
 \frac{H^x}{q-1}<\frac2{B^x}\le\frac1{128}<\frac1{32}.
\]
The same geometric-series argument now proves
$0\le r^0<q^k/32$ for each of the ten range digits. The common norm
therefore recovers every ordinary remainder and its intended quotient.
Restoring each deleted interval slack by $s_i=q^k-r_i^0-1$ gives
$s_i\ge0$; a fresh $G$-representation supplies the former product sign
certificate if it is wanted. The unchanged divisibility and prime-power
range-decoding proof applies.

For completeness, start with a finite refutation and use the parent's
cofinal horizon, beta-CRT construction, amplified CRT representatives,
even dominant semantic anchor, and private Pell witnesses. These genuine
witnesses satisfy both coefficient estimates just proved. Set
$v_i=M/L_i$ and $R_i=v_i r_i^0$. The three ratio equations and eleven
scaled digits hold. Every $R_i$ is below $M/32$, so
Theorem~\ref{thm:common} supplies the three norm witnesses. Replace the
four old denominator products by $NM$, retain all other semantic data,
and apply the supplied integral elimination and translation maps. This
produces a zero of the final polynomial from any finite refutation.

\section{The explicit quartic and its logical calibration}\label{sec:calibration}

\begin{theorem}[The released direct equation]\label{thm:main}
Let $Q_0,\ldots,Q_{181}\in\Z[z_0,\ldots,z_{251}]$ be the fixed
residuals printed in Appendix~\ref{app:residuals} and serialized in
\file{direct/equation/equation.json}. Define
\[
 P_{252}=\sum_{j=0}^{181}Q_j^2.
\]
Then $P_{252}$ has degree four, 5,340 collected monomials, and coefficient
height 128,446. Using the pinned direct-calculus calibration,
\[
 \ZFC\vdash
 (\exists\mathbf z\in\Z^{252}\;P_{252}(\mathbf z)=0)
 \longleftrightarrow
 (\exists p\;\operatorname{Refutation}_{\ZFC}(p)).
\]
Therefore ZFC proves unsolvability of this fixed equation equivalent to
$\Con(\ZFC)$.
\end{theorem}
\begin{proof}
An integer sum of squares is zero exactly when every summand is zero.
The literal Horn factoring, integral monic projections, and integer
translations preserve existential solvability, with explicit inverse
maps. The preceding sections prove the two directions between the new
arithmetic source and the parent's interpreted direct arithmetic source.
The latter decodes a finite refutation in the numerical calculus C222.
The retained endpoint interfaces are
\file{KEndpoint.c222_iff_zf_refutation} and
\file{ConstructibleChoice.zf_refutation_iff_zfc_refutation}.
Their composition gives the displayed refutation equivalence. Negating
it gives the consistency formulation. The degree and coefficient counts
follow by exact collection of the supplied residual squares.
\end{proof}

The theorem uses the repository's existing ordinary arithmetic and
logical calibration. The new contribution is the remainder representation
and its two-stage recovery proof, together with the exact polynomial
transformation. The archive preserves the relevant inherited proof texts,
source hashes, and calibration names. It does not claim that a Python
coefficient audit proves an unbounded arithmetic theorem, nor that the
whole transitive Lean/Rocq development was rebuilt in this iteration.

No consistency assumption is used to prove the equivalence. Under
consistency, the equation has no integer zero; if a syntactic contradiction
exists, the constructive direction produces a zero. Neither alternative
is decided by this report.

\subsection{Why precisely thirty variables disappear}

Eleven separate interval predicates previously needed three ternary sign
witnesses each, for 33 variables. The new construction adds three common
divisor ratios and three witnesses for the one norm predicate. This
saves 27. Replacing the four products $Nq^k$ by $NM$ saves three more.
The old interval slacks were already removable by affine elimination;
they are not counted again. The final literal count is consequently
$282-27-3=252$.

\section{A smaller denominator in the compact Pell carrier}\label{sec:carrier}

This secondary result improves the ten-positive-variable equations in the
incoming \emph{Certificate Reductions} package~\cite{certreductions}.
It is a separate K215/native route. Its constants are the complete fixed
indexed tuple of that package, not the direct quartic's small coefficients.

The predecessor's first Pell base is $P_0=2XY^2+1$, with
$X=wq$ and $Y=q^3s$. It uses a carrier denominator $D=Y^2$.
The following proposition identifies the actual denominator requirements.

\begin{proposition}[Carrier substitution]\label{prop:carrier}
Keep the predecessor's first Pell base and all its definitions other than
the carrier denominator. Any positive polynomial value $D$ satisfying
\[
 D>2,\qquad q^4\mid D,\qquad D\mid4XY^2
\]
on every positive input can replace $Y^2$ in the carrier without changing
its equivalence with the native refutation condition.
\end{proposition}
\begin{proof}
The unchanged product-square extraction makes the rational carrier an
integer and yields the same main and first Pell representations. Its sign
transport uses only $D>2$. It therefore gives the same positive outer
data and $R<q^4$. Clearing the integral carrier gives
$D\mid hHe$, with $h=q-1$ and $H\equiv-1\pmod q$.
Since $q^4\mid D$ and $\gcd(q,hH)=1$, we obtain $q^4\mid e$.

The first Pell congruence is $k\equiv2n_0\pmod{4XY^2}$, and the
unchanged growth argument gives $1\le n_0<R$. Thus
\[
 q^4\mid2n_0-R-1,\qquad |2n_0-R-1|<R<q^4,
\]
so $2n_0=R+1$. Returning to the full congruence gives
$4XY^2\mid e$, hence $D\mid e$. The rest of the projection and
native decoding is unchanged.

Conversely, canonical first Pell witnesses give $e$ a positive multiple
of $4XY^2$, so $e/D$ is a positive integer. The predecessor's fresh
strong and auxiliary witnesses make the padding coordinate
$n=gy+Wt+e/D-\kappa$ positive, because $g>\kappa$ and the other
summands have the stated nonnegative signs. The carrier equation holds.
\end{proof}

Set $D=XY$. Both divisibilities are immediate, and the evaluation can
compute $D=XY$, then $V_0=DY$, in place of $Y^2$, then $V_0=XY^2$.
The number of multiplication gates is unchanged. The denominator degree
falls from eight to six, lowering the final degree by four.

\begin{center}
\begin{tabular}{@{}lrrr@{}}
\toprule
Carrier & Positive variables & Binary operations & Degree\\
\midrule
Parent $D=Y^2$, weight $V^2$ & 10 & 78 & 630\\
New $D=XY$, weight $V^2$ & 10 & 78 & 626\\
Parent $D=Y^2$, weight $bV$ & 10 & 79 & 468\\
New $D=XY$, weight $bV$ & 10 & 79 & 464\\
New $D=q^4$, weight $bV$ & 10 & 80 & 460\\
\bottomrule
\end{tabular}
\end{center}

The full formulas, literal division-free DAGs, and the nonzero leading
homogeneous forms are in \file{compact/}. The complete predecessor proof
is retained. These are modest degree improvements, not a claim to a
new universal variable bound. The $79/622$ variant from $D=q^4$ and
weight $V^2$ is dominated by $79/464$ and is not promoted.

\section{Native coefficients: a completed controller and a smaller radix}\label{sec:native}

The direct quartic is the principal result of this report. A separate
native route keeps fifteen positive integer variables and a 66-operation
polynomial of degree 157, while reducing the size of its four fixed
integer coefficients. This section separates a completed prior proposal
from a new change to the encoding.

\subsection{The polynomial is fixed before the controller is chosen}

The fifteen independent positive variables are
\[
 J,F,\alpha,z,f,h,i,T,s,w,\tau,\eta,\zeta,y,\gamma.
\]
For a specified whole coefficient tuple $(\beta,K_0,M_C,M_F)$, define
the following polynomial abbreviations:
\[
\begin{aligned}
 q&=\beta J+1,& C&=q-F-\alpha,\\
 X&=wq,& Y&=sq^3,\\
 E&=XY,& k&=\eta+\zeta,\\
 c&=kY+\eta,& \ell&=EkY,\\
 a&=E+Y,& H&=4a+3,\\
 \Delta&=a^2+H,& D&=X+ac+\gamma H,\\
 V&=f^2(cT-R)-c,& S&=i\Delta c^2.
\end{aligned}
\]
\[
 R=[(q-1)(q-F)+\alpha](q^2-1)+(M_C+qM_F)J.
\]
Put
\begin{align*}
 N_0&=\tau^2-\ell(\ell+k),&N_1&=D^2-\Delta c^2,\\
 N_a&=S^2(V^2-y^2)+y^2,&N_k&=k-hE-R,\\
 N_t&=(K_0+w)C+(q-F)-z(q-1),&N_s&=\Delta f^2-S^2.
\end{align*}
The single equation is
\begin{equation}\label{eq:p66}
 \boxed{N_0N_1N_aN_kN_tN_s-\Delta=0.}
\end{equation}
Its exact straight-line circuit has 38 multiplications, 15 additions,
and 13 subtractions. The apparent top terms of $N_1$ cancel: with $Z=X+\gamma H$,
$N_1=2acZ+Z^2-Hc^2$. The six factor degrees are therefore
$22,18,62,7,2,46$. Setting all inputs equal to an indeterminate $t$
gives leading term $2^{17}\beta^{91}t^{157}$, so the degree is exactly
157. The four named constants are not quantified unknowns. The controller prime below is not part of the coefficient
tuple; it is an explicit witness used in the completeness proof.

\subsection{Closing the conditional radix-$2^{49}$ construction}

The repository's shared-marker radix49 proposal~\cite{radix49} already
contains the encoding and its ordinary native/Pell proof. Its remaining
conditions were finite controller conditions. We complete precisely
those conditions and materialize the entire coherent tuple.

Let $\rho=2^{49}$, $L=8807$, $B=\rho^L$, and $d=49L=431543$.
Choose the fixed prime
\[
 \lambda=434011>d,\qquad \lambda-1=1887\cdot230.
\]
For $H_{\rm ctrl}=1+BK_0+B^2$, exact computation gives:
\begin{center}
\begin{tabular}{@{}rrrr@{}}
\toprule
Modulus & $B$ & $K_0$ & $H_{\rm ctrl}$\\
\midrule
7 & 4 & 5 & 2\\
49 & 18 & 47 & 44\\
8807 & 1901 & 973 & 3135\\
434011 & 398487 & 335571 & 6522\\
\bottomrule
\end{tabular}
\end{center}
The exact orders are
\[
 \operatorname{ord}_{49}(B)=3,\quad
 \operatorname{ord}_{8807}(B)=629,\quad
 \operatorname{ord}_{434011}(B)=18870.
\]
The required valuation-one lift is
\[
 B^{\lambda-1}\equiv96356084144
 =1+222013\lambda\pmod{\lambda^2}.
\]
Since $0<222013<\lambda$, the lift has valuation exactly one.
The full modulus 49 is checked explicitly, rather than only its prime
divisor seven. Deterministic trial division proves the asserted primes;
annihilating powers and every maximal proper-power test prove the orders.

These facts discharge all the finite premises of the inherited
conditional theorem. Its carry, phase-alignment, safe-dummy, rank-tuning,
Pell, and refutation arguments then give the ZFC-provable equivalence
between positive unsolvability of~\eqref{eq:p66} at this tuple and
$\Con(\ZFC)$. The new work here is the finite completion and complete
coefficient specialization; the shared-marker architecture is prior work.

The materialized radix49 tuple has bit lengths
$(431543,433602,431543,431544)$ in the port order
$(\beta,K_0,M_C,M_F)$, totaling 216,030 raw unsigned bytes. Its four
deterministic gzip files total 8,774 bytes. The corresponding matched
radix53 and radix55 tuples occupy 233,666 and 242,484 raw bytes.
Raw coefficient storage, compressed storage, and printed-page length
are different metrics. The complete coefficient recipe, binary files,
standalone verifier, and inherited proof texts are supplied under
\file{radix49/}.

\subsection{Complemented labels permit a radix-$2^{47}$ encoding}

The radix49 completion suggests a further change to the coefficient
construction. The previous radix47 local obstruction~\cite{radix47obstruction}
concerned unchanged selector coefficients: its predecoded carry allowance
was too large to infer that a cell selected at most one tile. Instead of
discarding that obligation, change the numerical labels so the full
coefficient sum becomes smaller.

For each original tile $(w_i,n_i,p_i,s_i)$, apply the bijections
\[
 w_i'=255-w_i,\quad p_i'=255-p_i,\quad
 n_i'=511-n_i,\quad s_i'=511-s_i.
\]
The west/east labels use the same eight-bit bijection and the north/south
labels the same nine-bit bijection. Thus edge equality and inequality,
including the typed Start and forbidden-color conditions, are preserved.
The seven physical marker bits and the occupancy bits are kept. An empty
cell still has zero physical label bits; complementation is the relabeling
of genuine tiles, not an unconditional complement of every empty bit.

Put
\[
 u_i=w_i+2^8n_i+2^{17}p_i+2^{25}s_i.
\]
The complemented packed label is $2^{34}-1-u_i$, so the new selector
coefficient is
\begin{equation}\label{eq:complement}
 c_i=1+2^{10}(u_i+1)\qquad(0\le i\le942).
\end{equation}
Every nonselector coefficient and physical support position is as in the
shared-marker construction. Exact reconstruction of all 943 tiles gives
\[
 \sum_e c_e=3,953,059,820,717,999,
 \qquad\max_i c_i=15,049,678,914,561<2^{44}.
\]
The old complete coefficient sum was 12,917,846,596,842,415. Both sums
include the nonselector controls; the reduction is not obtained by
omitting any coefficient from the proof.

\subsubsection*{The full new coefficient tuple}

Now set $\rho=2^{47}$, $L=8807$, $B=\rho^L=2^{413929}$,
$\beta=B-1$, and $\Lambda=\rho^{4393}+\rho^{5495}$.
Besides~\eqref{eq:complement}, the complete nonzero controls are
\begin{align*}
 c_{1023+j}&=2^{10+j},&c_{1003+j}&=2^{27+j}&&(0\le j\le7),\\
 c_{1013+j}&=2^{18+j},&c_{983+j}&=2^{35+j}&&(0\le j\le8),\\
 c_e&=2^{44}&&&
 e&\in\{992,1011,1022,1031,1095,2195,3296\}.
\end{align*}
All other $c_e$ vanish. Let $\mathcal E$ be their support together with
the two coefficient-zero dummy positions 1063 and 1064. It has 986
positions. Define
\begin{align*}
 K_0={}&\rho^{30}+\rho^{62}+\rho^{1100}+\rho^{2201}+\rho^{8849}
       +\sum_{e=0}^{3296}c_e(\rho^{4393-e}+\rho^{5495-e}),\\
 M_C={}&\beta-\sum_{e\in\mathcal E\setminus\{1\}}\rho^e-2\rho^{1063},\\
 C_{\rm core}={}&\sum_{e=1013}^{1022}\rho^e+
              \sum_{e=1065}^{1073}\rho^e+\rho^{2195}+\rho^{3296},\\
 M_F={}&2B+3+(2^{47}-2)\Lambda+C_{\rm core}-\rho^{8792}
                +(2^{28}-3)\rho^{7694}-\rho^{7691}.
\end{align*}
These are fixed integers. Every exponent and finite sum limit is fixed.
The exact 943-tile data and the complete 3297-slot coefficient array are
supplied; all four binary integers are also materialized. In particular,
no unspecified tile enumeration or proof search is needed to define them.

The positive native high-mask padding behind the compressed $M_F$ formula
is fifteen full fields at 8792--8806, three full fields at 7691--7693,
and bits 2--27 at field 7694. Its population is
$15\cdot47+3\cdot47+26=872$. Together with the unchanged 114 core
tests, this gives 986 ones, exactly matching the physical low support.
The source shift by $\beta$ is applied once. All four ports are changed
coherently; none is imported numerically from the radix49 tuple.

\subsubsection*{Why the old carry obstruction no longer applies}

The complete predecoded convolution bound, including the exceptional
dummy and every temporal phase, maps the carry interval $0,\ldots,29$
into itself. Thus
\[
 \#\{\hbox{selected tiles}\}+k\le943+29=972<1024.
\]
The ten low tested bits therefore recover zero or one selected tile.
The remaining payload bits recover the complemented label and, by the
same shared-marker counting argument, all seven marker bits individually.

After this decoding, the selector maximum is below $\rho/8$ and the
entire nonselector control sum is $\rho-1024$. The complete decoded
convolution and temporal estimate, even with the earlier incoming carry,
is strictly below $13\rho$. Hence the actual outgoing carries are at
most 12. This bound protects the anchor predecessors and all subsequent
local resets.

At radix47 an occupied clause has raw value $\rho+1$. Its canonical
digit is one and it emits carry one. This carry must not be silently
dropped. The proof explicitly resets it before every later tested copy,
anchor, padding, or origin field. For example, the larger clause
predecessor has raw bound at most $1+7\cdot2^{44}-2^{10}$, and adding
12 still leaves it below $\rho$. The two phase-anchor predecessor
identities also leave room for every temporal phase plus carry 12.
They recover integral temporal phase and whole-cell alignment as before.
The raw padding predecessors at 7690 and 8791 are bounded by a single
selector coefficient plus the incoming carry, so both reset. The complete
zero fields 7691--7695 and 8792--8806 then validate the new padding.

These arguments preserve the positive-source slack, all simultaneous
safe-dummy changes, and the actual rank identity. They supply the same
canonical $R\equiv3\pmod4$ and
$\pc(R)=3dN+2$ interface required by the P66 converse. The full ordinary
proof and the independent review are included in \file{radix47/}.

\subsubsection*{The new finite controller}

Here $d=47\cdot8807=413929$. Choose
\[
 \lambda=1012691>d,\qquad
 \lambda-1=10\cdot101269,
 \qquad101269=\operatorname{lcm}(23,4403).
\]
The exact residues are:
\begin{center}
\begin{tabular}{@{}rrrr@{}}
\toprule
Modulus & $B$ & $K_0$ & $1+BK_0+B^2$\\
\midrule
47 & 12 & 5 & 17\\
8807 & 2677 & 4079 & 5042\\
1012691 & 959313 & 1005855 & 834350\\
\bottomrule
\end{tabular}
\end{center}
The exact orders at these three moduli are respectively
$23$, $4403$, and $1012690$. Finally,
\[
 B^{\lambda-1}\equiv1+282338\lambda\pmod{\lambda^2}.
\]
All controller units are nonzero and the lift has valuation one. The
same cofinal CRT and Boolean-subset rank-tuning proof now applies with
the new $d$, prime, spacing, and tuple. Together with the complete
carry/alignment proof, it establishes the ZFC-provable equivalence
between positive unsolvability of~\eqref{eq:p66} at this tuple and
$\Con(\ZFC)$.

\begin{center}
\begin{tabular}{@{}lrrrrr@{}}
\toprule
Tuple & $\beta$ bits & $K_0$ bits & $M_C$ bits & $M_F$ bits & Raw bytes\\
\midrule
Radix55 & 484385 & 486696 & 484385 & 484386 & 242484\\
Radix53 & 466771 & 468998 & 466771 & 466772 & 233666\\
Completed radix49 & 431543 & 433602 & 431543 & 431544 & 216030\\
New radix47 & 413929 & 415904 & 413929 & 413930 & \textbf{207214}\\
\bottomrule
\end{tabular}
\end{center}

The new tuple saves 26,452 raw bytes, about 11.32\%, against the verified
radix53 tuple. Its compressed four-coefficient payload is 8,356 bytes.
The scalar $B=2^{413929}$ is also smaller than the collection's previously
stated $2^{429959}$ native-base endpoint; this statement concerns that
scalar, not a comparison of every metric of the two different encodings.

Label complementation itself has predecessors in the repository. The new
result is this complete indexed943/shared-marker/radix47 composition,
the corrected carry bounds and resets, its entire coherent tuple, and the
completed controller. It is not a claim to have invented complementation
or to have proved a globally smallest coefficient tuple.



\section{Verification, reproducibility, and limits}\label{sec:verification}

The final archive includes complete defining objects, not just programs
that might eventually find their coefficients. The direct equation can
be read directly from the residual JSON or the printed appendix. Its
expanded JSON contains the whole collected quartic. No search for a
solution or consistency-dependent computation is part of coefficient
generation.

The direct finite audit checks the exact deletion and replacement of
source rows, all remaining guards, every ratio equation, every scaled
digit, the shared-product aliases, all inherited Horn factoring edits,
205 monic eliminations with restoration polynomials, and the final
integer translation. It independently recollects every coefficient of
the 5,340-term polynomial. Deliberately altered ratio, digit, norm,
restoration, translation, and expanded-coefficient data are rejected.
Normal Python and optimized Python both retain the checks.

The finite native audits independently reconstruct complete coefficient
tuples, check prime/order/lift certificates and controller units, and
verify the literal P66 circuit and its degree. Exact integer arithmetic
is used throughout. These checks supplement the ordinary carry,
Pell, CRT, and refutation proofs; they do not replace them.

For the direct component, run the following commands from the archive root:
\begin{verbatim}
python3 -B direct/code/verify.py
python3 -B -O direct/code/verify.py
python3 -B direct/code/verify.py --rebuild
\end{verbatim}
The optional rebuild uses a temporary directory and compares nineteen
generated JSON objects byte for byte. It requires no network or repository
checkout. All arithmetic checkers use only the Python standard library.
The enclosing \file{README.md} gives the exact native and compact commands
and the report's LaTeX compilation command. Source licensing and AI
assistance notices are retained in each component.


The result is best described as a concrete improvement in several
explicit constructions. It is not established to be the globally
simplest Diophantine equation for consistency. In particular, fewer
variables alone would favor some different higher-degree constructions,
while smaller coefficients favor the direct quartic. The archive keeps
those comparisons separate rather than counting one measure as all of
them.

\section{Further research questions}

\begin{enumerate}
\item \textbf{A systematic common-remainder compiler.} Which digit-based
      Diophantine constructions have canonical remainders uniformly much
      smaller than their moduli? Theorem~\ref{thm:common} applies to any
      such family, but proving the margin before decoding is the main
      semantic obligation. The divisors need only share an explicitly
      certified common multiple.
\item \textbf{A better common multiple.} The choice $M=q^x$ is convenient
      because it already exists. A different grouping or smaller common
      multiple could reduce support or the number of ratio coordinates.
      Any saving must include the variables needed to certify the new
      multiple and its exact quotients.
\item \textbf{The three initial sign guards.} The remaining arithmetic
      bootstrap still uses the initial remainder, slack, and output-slack
      signs. Replacing them requires a new noncircular proof of $B>0$
      and $q\ge2$. The common-remainder lemma alone does not justify
      deleting these guards.
\item \textbf{Degree and support after coordinate changes.} The current
      monic eliminations optimize variable count subject to a sparse
      substitution cap. Alternative integral coordinates could lower the
      5,340-term support or coefficient height while preserving 252
      variables. Such improvements should be measured on collected
      polynomials and accompanied by inverse maps.
\item \textbf{Native label choices and carry proofs.} Label names are
      semantically arbitrary, but their binary values affect selector
      coefficients and carry bounds. A finite optimization over valid
      color bijections is potentially useful. It must preserve complete
      boundary and occupancy semantics, and must prove all downstream
      resets, including any newly created clause carry. Below radix47, the
      current uniform clause-predecessor bound $7\cdot2^{44}-2^{10}$
      itself exceeds the next radix; a smaller global coefficient sum
      alone does not resolve that local obstruction.
\item \textbf{Formalizing the new local theorem.} Lemma~\ref{lem:rescale},
      the common norm theorem, and Lemma~\ref{lem:binom} are small,
      independent formalization targets. Connecting their verified source
      transformation to the complete literal polynomial would provide a
      focused proof-assistant contribution without claiming to rebuild
      the entire inherited compiler at once.
\end{enumerate}

\begin{thebibliography}{99}
\bibitem{repository} V. Reshetnikov, \emph{ProveIt}, source snapshot
\sourcepin, 25 October 2026.
\url{https://github.com/VladimirReshetnikov/ProveIt/tree/55d1f72aedebf892e446c90d35fa1ae37689b0f2/Computability/DiophantineConsistency}.
\bibitem{round17} ProveIt incoming research package,
\emph{Cofinal Anchors and Height One, Round 17}, 25 October 2026,
particularly \file{direct/STRICT_BUDGET_PROOF.md},
\file{direct/ORIENTATION_FIRST_PROOF.md}, and the
\file{strict_dimension_centered} literal equation. The exact source
archive is \file{docs/incoming/ConZFC_Cofinal_Anchors_Height_One_Round17_2026-10-25.zip}
at~\cite{repository}. Relevant proof texts and exact input objects are
retained in this release.
\bibitem{certreductions} ProveIt incoming research package,
\emph{Certificate Reductions}, 25 October 2026,
\file{compact/PROOF.md} and its two fixed 78/630 and 79/468 circuits.
Source archive: \file{docs/incoming/ConZFC_Certificate_Reductions_2026-10-25.zip}
at~\cite{repository}.
\bibitem{radix49} ProveIt, \emph{Seven shared marker weights and a
conditional radix49 adapter},
\file{research/compact66-shared-marker-radix49/README.md} at~\cite{repository}.
The architecture and its ordinary proof precede the finite certificate
completion in this report.
\bibitem{radix47obstruction} ProveIt, \emph{A radix47 alias of the
isolated shared-marker clause test},
\file{research/compact66-shared-marker-radix47-clause-alias/README.md}
at~\cite{repository}. This is a scoped obstruction to a former local
argument, not an impossibility theorem for radix47 encodings.
\bibitem{matiyasevich} Y. Matiyasevich, \emph{Hilbert's Tenth Problem},
MIT Press, 1993; author's online preface and commentary:
\url{https://logic.pdmi.ras.ru/~yumat/H10Pbook/preface.htm}.
\bibitem{jones} J. P. Jones, Universal Diophantine equation,
\emph{Journal of Symbolic Logic} \textbf{47} (1982), 549--571.
\url{https://doi.org/10.2307/2273588}.
\bibitem{bayer} J. Bayer, M. David, M. Hassler, Y. Matiyasevich, and
D. Schleicher, \emph{Diophantine Equations over $\Z$: Universal Bounds
and Parallel Formalization}, 2025.
\url{https://arxiv.org/abs/2506.20909}.
Related checked development:
\url{https://isa-afp.org/entries/Diophantine_Universal_Pairs.html}.
\end{thebibliography}

\appendix
\section{The complete direct polynomial}\label{app:residuals}

All $z_0,\ldots,z_{251}$ are unrestricted integers. The following is the
complete list of fixed residuals. No abbreviation in this appendix is
an additional quantified variable. The requested single equation is
\[
 \boxed{\displaystyle\sum_{j=0}^{181}Q_j(\mathbf z)^2=0.}
\]
The authoritative variable order, residual coefficients, and full
collected expansion are also supplied in machine-readable form. The
appendix is generated directly from that final JSON, and its metadata
and coefficient expansion are independently checked.

\begingroup
\small
\setlength{\jot}{3pt}
\begin{align*}
Q_{0} &={}-z_{3}^{2} + z_{4}\\[4pt]
Q_{1} &={}z_{0} - z_{0}z_{2} - z_{2} + z_{9}\\[4pt]
Q_{2} &={}-1 + z_{1} - z_{2}z_{7} - z_{3}\\[4pt]
Q_{3} &={}-1 + z_{1} - z_{8}z_{9}\\[4pt]
Q_{4} &={}1 + z_{4} - z_{5}z_{20} + z_{18}\\[4pt]
Q_{5} &={}-1 + 2z_{10} - z_{10}^{2} - 2z_{10}z_{11} + 2z_{11} - z_{11}^{2} + z_{21}\\[4pt]
Q_{6} &={}-5 - 30z_{10} - 36z_{10}^{2} - 96z_{10}z_{11} - 24z_{10}z_{21} - 40z_{11}\\*
&\quad{} - 64z_{11}^{2} - 32z_{11}z_{21} - 10z_{21} - 4z_{21}^{2} + z_{22}\\[4pt]
Q_{7} &={}-2 + 2z_{10} - z_{10}^{2} - 4z_{10}z_{22} + 4z_{22} - 4z_{22}^{2} + z_{23}\\[4pt]
Q_{8} &={}1 + 6z_{10} - 36z_{10}^{2} - 192z_{10}z_{22} - 24z_{10}z_{23} + 16z_{22}\\*
&\quad{} - 256z_{22}^{2} - 64z_{22}z_{23} + 2z_{23} - 4z_{23}^{2} + z_{24}\\[4pt]
Q_{9} &={}-2 - z_{11}^{2} - 2z_{11}z_{12} - z_{12}^{2} + z_{25}\\[4pt]
Q_{10} &={}-9 + 7z_{10} - z_{10}^{2} - 4z_{10}z_{11} - 4z_{10}z_{25} + 12z_{11}\\*
&\quad{} - 4z_{11}^{2} - 8z_{11}z_{25} + 12z_{25} - 4z_{25}^{2} + z_{26}\\[4pt]
Q_{11} &={}-1 + 2z_{10} - z_{10}^{2} - 2z_{10}z_{12} + 2z_{12} - z_{12}^{2} + z_{27}\\[4pt]
Q_{12} &={}-6z_{10} - 144z_{10}^{2} - 192z_{10}z_{11} - 192z_{10}z_{12}\\*
&\quad{} - 48z_{10}z_{21} - 48z_{10}z_{27} - 8z_{11} - 64z_{11}^{2} - 128z_{11}z_{12}\\*
&\quad{} - 32z_{11}z_{21} - 32z_{11}z_{27} - 64z_{12}^{2} - 32z_{12}z_{21}\\*
&\quad{} - 32z_{12}z_{27} - 2z_{21} - 4z_{21}^{2} - 8z_{21}z_{27} - 4z_{27}^{2} + z_{28}\\[4pt]
Q_{13} &={}-4z_{26}^{2} - 8z_{26}z_{28} - 4z_{28}^{2} + z_{29}\\[4pt]
Q_{14} &={}-6 + 16z_{12} - 16z_{12}^{2} - 16z_{12}z_{25} - 16z_{12}z_{27} + 8z_{25}\\*
&\quad{} - 4z_{25}^{2} - 8z_{25}z_{27} + 8z_{27} - 4z_{27}^{2} + z_{30}\\[4pt]
Q_{15} &={}1 + 6z_{10} - 36z_{10}^{2} - 96z_{10}z_{11} - 48z_{10}z_{12}\\*
&\quad{} - 24z_{10}z_{21} - 48z_{10}z_{27} - 24z_{10}z_{30} + 8z_{11} - 64z_{11}^{2}\\*
&\quad{} - 64z_{11}z_{12} - 32z_{11}z_{21} - 64z_{11}z_{27} - 32z_{11}z_{30}\\*
&\quad{} + 8z_{12} - 16z_{12}^{2} - 16z_{12}z_{21} - 32z_{12}z_{27} - 16z_{12}z_{30}\\*
&\quad{} + 2z_{21} - 4z_{21}^{2} - 16z_{21}z_{27} - 8z_{21}z_{30} + 8z_{27}\\*
&\quad{} - 16z_{27}^{2} - 16z_{27}z_{30} + 4z_{30} - 4z_{30}^{2} + z_{31}\\[4pt]
Q_{16} &={}-10 - 5z_{10} - z_{10}^{2} - 12z_{10}z_{11} - 16z_{10}z_{12}\\*
&\quad{} - 4z_{10}z_{25} - 36z_{11} - 36z_{11}^{2} - 96z_{11}z_{12} - 24z_{11}z_{25}\\*
&\quad{} - 48z_{12} - 64z_{12}^{2} - 32z_{12}z_{25} - 12z_{25} - 4z_{25}^{2} + z_{32}\\[4pt]
Q_{17} &={}-2 - 24z_{10} - 36z_{10}^{2} - 12z_{10}z_{11} - 96z_{10}z_{12}\\*
&\quad{} - 24z_{10}z_{27} - 3z_{11} - z_{11}^{2} - 16z_{11}z_{12} - 4z_{11}z_{27}\\*
&\quad{} - 32z_{12} - 64z_{12}^{2} - 32z_{12}z_{27} - 8z_{27} - 4z_{27}^{2} + z_{33}\\[4pt]
Q_{18} &={}2 - 4z_{32}^{2} - 8z_{32}z_{33} - 4z_{33}^{2} + z_{34}\\[4pt]
Q_{19} &={}-12 + 16z_{10} - 4z_{10}^{2} - 8z_{10}z_{12} - 8z_{10}z_{21}\\*
&\quad{} - 8z_{10}z_{25} + 16z_{12} - 4z_{12}^{2} - 8z_{12}z_{21} - 8z_{12}z_{25}\\*
&\quad{} + 16z_{21} - 4z_{21}^{2} - 8z_{21}z_{25} + 16z_{25} - 4z_{25}^{2} + z_{35}\\[4pt]
Q_{20} &={}-3 - 24z_{10} - 64z_{10}^{2} - 160z_{10}z_{12} - 64z_{10}z_{25}\\*
&\quad{} - 32z_{10}z_{27} - 32z_{10}z_{35} - 26z_{12} - 100z_{12}^{2}\\*
&\quad{} - 80z_{12}z_{25} - 40z_{12}z_{27} - 40z_{12}z_{35} - 8z_{25} - 16z_{25}^{2}\\*
&\quad{} - 16z_{25}z_{27} - 16z_{25}z_{35} - 6z_{27} - 4z_{27}^{2} - 8z_{27}z_{35}\\*
&\quad{} - 4z_{35} - 4z_{35}^{2} + z_{36}\\[4pt]
Q_{21} &={}-1 - 18z_{10} - 36z_{10}^{2} - 96z_{10}z_{11} - 24z_{10}z_{21} - 24z_{11}\\*
&\quad{} - 64z_{11}^{2} - 32z_{11}z_{21} - 6z_{21} - 4z_{21}^{2} + z_{37}\\[4pt]
Q_{22} &={}-2 + 2z_{10} - z_{10}^{2} - 4z_{10}z_{37} + 4z_{37} - 4z_{37}^{2} + z_{38}\\[4pt]
Q_{23} &={}1 + 6z_{10} - 36z_{10}^{2} - 192z_{10}z_{37} - 24z_{10}z_{38} + 16z_{37}\\*
&\quad{} - 256z_{37}^{2} - 64z_{37}z_{38} + 2z_{38} - 4z_{38}^{2} + z_{39}\\[4pt]
Q_{24} &={}-z_{11}^{2} - 4z_{11}z_{39} - 4z_{39}^{2} + z_{40}\\[4pt]
Q_{25} &={}-6 - 30z_{10} - 64z_{10}^{2} - 128z_{10}z_{11} - 64z_{10}z_{21} - 32z_{11}\\*
&\quad{} - 64z_{11}^{2} - 64z_{11}z_{21} - 14z_{21} - 16z_{21}^{2} + z_{41}\\[4pt]
Q_{26} &={}-4 + 8z_{11} - 64z_{11}^{2} - 96z_{11}z_{12} - 32z_{11}z_{25}\\*
&\quad{} - 32z_{11}z_{41} + 6z_{12} - 36z_{12}^{2} - 24z_{12}z_{25} - 24z_{12}z_{41}\\*
&\quad{} + 2z_{25} - 4z_{25}^{2} - 8z_{25}z_{41} + 4z_{41} - 4z_{41}^{2} + z_{42}\\[4pt]
Q_{27} &={}8z_{10} - 64z_{10}^{2} - 96z_{10}z_{12} - 32z_{10}z_{27} - 32z_{10}z_{42}\\*
&\quad{} + 6z_{12} - 36z_{12}^{2} - 24z_{12}z_{27} - 24z_{12}z_{42} + 2z_{27}\\*
&\quad{} - 4z_{27}^{2} - 8z_{27}z_{42} + 4z_{42} - 4z_{42}^{2} + z_{43}\\[4pt]
Q_{28} &={}-z_{12}^{2} - 2z_{12}z_{13} - z_{13}^{2} + z_{44}\\[4pt]
Q_{29} &={}2z_{10} - z_{10}^{2} - 4z_{10}z_{12} - 4z_{10}z_{44} + 4z_{12} - 4z_{12}^{2}\\*
&\quad{} - 8z_{12}z_{44} + 4z_{44} - 4z_{44}^{2} + z_{45}\\[4pt]
Q_{30} &={}-35 + 24z_{10} - 4z_{10}^{2} - 8z_{10}z_{11} - 8z_{10}z_{25}\\*
&\quad{} - 8z_{10}z_{27} + 24z_{11} - 4z_{11}^{2} - 8z_{11}z_{25} - 8z_{11}z_{27}\\*
&\quad{} + 24z_{25} - 4z_{25}^{2} - 8z_{25}z_{27} + 24z_{27} - 4z_{27}^{2} + z_{46}\\[4pt]
Q_{31} &={}-16z_{10}^{2} - 32z_{10}z_{12} - 32z_{10}z_{27} - 32z_{10}z_{44}\\*
&\quad{} - 16z_{10}z_{45} - 16z_{10}z_{46} - 4z_{12} - 16z_{12}^{2} - 32z_{12}z_{27}\\*
&\quad{} - 32z_{12}z_{44} - 16z_{12}z_{45} - 16z_{12}z_{46} - 16z_{27}^{2}\\*
&\quad{} - 32z_{27}z_{44} - 16z_{27}z_{45} - 16z_{27}z_{46} - 4z_{44} - 16z_{44}^{2}\\*
&\quad{} - 16z_{44}z_{45} - 16z_{44}z_{46} - 2z_{45} - 4z_{45}^{2} - 8z_{45}z_{46}\\*
&\quad{} - 4z_{46}^{2} + z_{47}\\[4pt]
Q_{32} &={}z_{12} - 25z_{12}^{2} - 40z_{12}z_{27} - 20z_{12}z_{30} - 16z_{27}^{2}\\*
&\quad{} - 16z_{27}z_{30} - 4z_{30}^{2} + z_{48}\\[4pt]
Q_{33} &={}-1 + 8z_{11} - 4z_{11}^{2} - 8z_{11}z_{25} + 8z_{25} - 4z_{25}^{2} + z_{49}\\[4pt]
Q_{34} &={}-1 - 16z_{11}^{2} - 8z_{11}z_{12} - 32z_{11}z_{25} - 16z_{11}z_{49}\\*
&\quad{} - z_{12}^{2} - 8z_{12}z_{25} - 4z_{12}z_{49} - 16z_{25}^{2} - 16z_{25}z_{49}\\*
&\quad{} - 4z_{49}^{2} + z_{50}\\[4pt]
Q_{35} &={}1 + z_{11} - z_{11}^{2} - 4z_{11}z_{12} - 4z_{11}z_{50} - 4z_{12}^{2}\\*
&\quad{} - 8z_{12}z_{50} - 4z_{50}^{2} + z_{51}\\[4pt]
Q_{36} &={}-1 - 4z_{51} - 4z_{51}^{2} + z_{52}\\[4pt]
Q_{37} &={}-1 - z_{12}^{2} - 8z_{12}z_{51} - 4z_{12}z_{52} - 16z_{51}^{2}\\*
&\quad{} - 16z_{51}z_{52} - 4z_{52}^{2} + z_{53}\\[4pt]
Q_{38} &={}2 + z_{11} - z_{11}^{2} - 4z_{11}z_{12} - 4z_{11}z_{53} - 4z_{12}^{2}\\*
&\quad{} - 8z_{12}z_{53} - 4z_{53}^{2} + z_{54}\\[4pt]
Q_{39} &={}-4z_{48}^{2} - 8z_{48}z_{54} - 4z_{54}^{2} + z_{55}\\[4pt]
Q_{40} &={}-9 + 7z_{10} - z_{10}^{2} - 8z_{10}z_{48} - 4z_{10}z_{55} + 24z_{48}\\*
&\quad{} - 16z_{48}^{2} - 16z_{48}z_{55} + 12z_{55} - 4z_{55}^{2} + z_{56}\\[4pt]
Q_{41} &={}-1 - 4z_{56} - 4z_{56}^{2} + z_{57}\\[4pt]
Q_{42} &={}-z_{12}^{2} - 8z_{12}z_{56} - 4z_{12}z_{57} - 16z_{56}^{2} - 16z_{56}z_{57}\\*
&\quad{} - 4z_{57}^{2} + z_{58}\\[4pt]
Q_{43} &={}-z_{10} - z_{10}^{2} - 4z_{10}z_{12} - 4z_{10}z_{58} - 4z_{12} - 4z_{12}^{2}\\*
&\quad{} - 8z_{12}z_{58} - 4z_{58} - 4z_{58}^{2} + z_{59}\\[4pt]
Q_{44} &={}-10 - 12z_{59} - 4z_{59}^{2} + z_{60}\\[4pt]
Q_{45} &={}1 + 4z_{11} - 4z_{11}^{2} - 8z_{11}z_{25} - 16z_{11}z_{59} - 8z_{11}z_{60}\\*
&\quad{} + 4z_{25} - 4z_{25}^{2} - 16z_{25}z_{59} - 8z_{25}z_{60} + 8z_{59}\\*
&\quad{} - 16z_{59}^{2} - 16z_{59}z_{60} + 4z_{60} - 4z_{60}^{2} + z_{61}\\[4pt]
Q_{46} &={}z_{11} - 25z_{11}^{2} - 40z_{11}z_{25} - 20z_{11}z_{61} - 16z_{25}^{2}\\*
&\quad{} - 16z_{25}z_{61} - 4z_{61}^{2} + z_{62}\\[4pt]
Q_{47} &={}-10 - 12z_{54} - 4z_{54}^{2} + z_{63}\\[4pt]
Q_{48} &={}1 - 16z_{54}^{2} - 16z_{54}z_{62} - 16z_{54}z_{63} - 4z_{62}^{2}\\*
&\quad{} - 8z_{62}z_{63} - 4z_{63}^{2} + z_{64}\\[4pt]
Q_{49} &={}z_{12} - z_{12}^{2} - 8z_{12}z_{62} - 4z_{12}z_{64} - 16z_{62}^{2}\\*
&\quad{} - 16z_{62}z_{64} - 4z_{64}^{2} + z_{65}\\[4pt]
Q_{50} &={}-4z_{10} - 4z_{10}^{2} - 4z_{10}z_{12} - 8z_{10}z_{21} - z_{12} - z_{12}^{2}\\*
&\quad{} - 4z_{12}z_{21} - 4z_{21} - 4z_{21}^{2} + z_{66}\\[4pt]
Q_{51} &={}1 + 4z_{10} - 4z_{10}^{2} - 8z_{10}z_{21} - 8z_{10}z_{66} + 4z_{21}\\*
&\quad{} - 4z_{21}^{2} - 8z_{21}z_{66} + 4z_{66} - 4z_{66}^{2} + z_{67}\\[4pt]
Q_{52} &={}-2 - 4z_{10} - 16z_{10}^{2} - 32z_{10}z_{21} - 16z_{10}z_{65}\\*
&\quad{} - 16z_{10}z_{67} - 4z_{21} - 16z_{21}^{2} - 16z_{21}z_{65} - 16z_{21}z_{67}\\*
&\quad{} - 4z_{65}^{2} - 8z_{65}z_{67} - 2z_{67} - 4z_{67}^{2} + z_{68}\\[4pt]
Q_{53} &={}-2 - 4z_{68} - 4z_{68}^{2} + z_{69}\\[4pt]
Q_{54} &={}-z_{3}z_{15} + z_{70}\\[4pt]
Q_{55} &={}z_{14} - z_{15} - 2z_{24}z_{70} - z_{71}\\[4pt]
Q_{56} &={}z_{14} - z_{15} - 4z_{26}z_{70} - 2z_{29}z_{70} - z_{72}\\[4pt]
Q_{57} &={}z_{14} - z_{15} - 2z_{31}z_{70} - 6z_{70} - z_{73}\\[4pt]
Q_{58} &={}z_{14} - z_{15} - 4z_{32}z_{70} - 2z_{34}z_{70} - z_{74}\\[4pt]
Q_{59} &={}z_{14} - z_{15} - 2z_{36}z_{70} + 6z_{70} - z_{75}\\[4pt]
Q_{60} &={}z_{14} - z_{15} - 4z_{39}z_{70} - 2z_{40}z_{70} - z_{76}\\[4pt]
Q_{61} &={}z_{14} - z_{15} - 2z_{43}z_{70} - z_{77}\\[4pt]
Q_{62} &={}z_{14} - z_{15} - 2z_{47}z_{70} - 4z_{70} - z_{78}\\[4pt]
Q_{63} &={}-z_{10}z_{11} + z_{10}z_{12} + z_{11} + z_{11}z_{12} - z_{12}\\*
&\quad{} - z_{12}^{2} + z_{79}\\[4pt]
Q_{64} &={}1 + 2z_{13} + z_{13}^{2} - z_{79}^{2} + z_{80}\\[4pt]
Q_{65} &={}-z_{80}^{2} + z_{81}\\[4pt]
Q_{66} &={}-4z_{68}z_{70} - 2z_{69}z_{70} + 2z_{70} + z_{82}\\[4pt]
Q_{67} &={}-z_{14}^{2} + 2z_{14}z_{15} + 2z_{14}z_{82} - z_{15}^{2} - 2z_{15}z_{82}\\*
&\quad{} - z_{82}^{2} + z_{83}\\[4pt]
Q_{68} &={}-1 - z_{3}z_{11} + z_{84}\\[4pt]
Q_{69} &={}-z_{15}z_{84} + z_{86}z_{87} + z_{88}\\[4pt]
Q_{70} &={}-1 + z_{3} - z_{3}z_{10} + z_{85}\\[4pt]
Q_{71} &={}-1 + 2z_{3} - 2z_{3}z_{10} - 2z_{3}z_{21} + z_{86}\\[4pt]
Q_{72} &={}-z_{14}z_{85} + z_{87}\\[4pt]
Q_{73} &={}-1 - 6z_{3}z_{10} - 8z_{3}z_{11} - 2z_{3}z_{21} + z_{89}\\[4pt]
Q_{74} &={}z_{14}z_{84} - z_{15}z_{89} + z_{90}\\[4pt]
Q_{75} &={}-z_{14}z_{71} + z_{15}z_{71} + z_{91}\\[4pt]
Q_{76} &={}-z_{72}z_{91} + z_{92}\\[4pt]
Q_{77} &={}-z_{73}z_{92} + z_{93}\\[4pt]
Q_{78} &={}-z_{74}z_{93} + z_{94}\\[4pt]
Q_{79} &={}-z_{75}z_{94} + z_{95}\\[4pt]
Q_{80} &={}-z_{76}z_{95} + z_{96}\\[4pt]
Q_{81} &={}-z_{77}z_{96} + z_{97}\\[4pt]
Q_{82} &={}-z_{78}z_{97} + z_{98}\\[4pt]
Q_{83} &={}-z_{88}z_{98} + z_{99}\\[4pt]
Q_{84} &={}-z_{90}z_{99} + z_{100}\\[4pt]
Q_{85} &={}-z_{100}^{2} + z_{101}\\[4pt]
Q_{86} &={}-z_{2} - z_{2}z_{18} + z_{18} + z_{102}\\[4pt]
Q_{87} &={}-z_{16}z_{102} + z_{103}\\[4pt]
Q_{88} &={}-2z_{1}^{2} + 2z_{1}z_{14} + 2z_{1}z_{15} + 2z_{1}z_{103} + 2z_{1}z_{104}\\*
&\quad{} - z_{14}^{2} - 2z_{14}z_{103} - z_{15}^{2} - 2z_{15}z_{104} + z_{19}z_{106}\\*
&\quad{} - z_{81}z_{101} - z_{83}z_{101} - z_{103}^{2} - z_{104}^{2}\\[4pt]
Q_{89} &={}-z_{2}z_{17} + z_{17} - z_{17}z_{102} + z_{104}\\[4pt]
Q_{90} &={}-9z_{10} - 9z_{11} - 4z_{21} - 2z_{24} - 4z_{26} - 2z_{29} - 2z_{31}\\*
&\quad{} - 4z_{32} - 2z_{34} - 2z_{36} - 4z_{39} - 2z_{40} - 2z_{43} - 2z_{47}\\*
&\quad{} - 4z_{68} - 2z_{69} + z_{105}\\[4pt]
Q_{91} &={}-z_{107}z_{110} + z_{110} + z_{112}\\[4pt]
Q_{92} &={}-z_{5}z_{134} + z_{108}z_{234} - z_{111} - z_{112}z_{134}\\[4pt]
Q_{93} &={}z_{113}z_{235} - z_{115}z_{137} - z_{116}\\[4pt]
Q_{94} &={}z_{117}z_{235} - z_{118}z_{137} - z_{119}\\[4pt]
Q_{95} &={}z_{120}z_{235} - z_{121}z_{137} - z_{122}\\[4pt]
Q_{96} &={}z_{123}z_{235} - z_{124}z_{137} - z_{125}\\[4pt]
Q_{97} &={}z_{126}z_{236} - z_{128}z_{137} - z_{129}\\[4pt]
Q_{98} &={}z_{130}z_{236} - z_{131}z_{137} - z_{132}\\[4pt]
Q_{99} &={}-z_{19}z_{134} + z_{137}\\[4pt]
Q_{100} &={}z_{133} - z_{135}z_{137} - z_{136}\\[4pt]
Q_{101} &={}-z_{137}z_{139} + z_{138} - z_{140}\\[4pt]
Q_{102} &={}-z_{137}z_{142} + z_{141}z_{234} - z_{143}\\[4pt]
Q_{103} &={}-z_{137}z_{145} + z_{144}z_{235} - z_{146}\\[4pt]
Q_{104} &={}z_{147} - 2z_{147}z_{149} + z_{150}\\[4pt]
Q_{105} &={}-z_{147}z_{150} - z_{150} + z_{151}\\[4pt]
Q_{106} &={}2z_{148} + z_{148}^{2} + z_{150}^{2} - z_{151}^{2}\\[4pt]
Q_{107} &={}-z_{148}z_{153} - z_{153} + z_{154}\\[4pt]
Q_{108} &={}2z_{152} + z_{152}^{2} + z_{153}^{2} - z_{154}^{2}\\[4pt]
Q_{109} &={}z_{4}^{2} - 2z_{4}z_{148} - 2z_{148} + z_{155}\\[4pt]
Q_{110} &={}-1 - z_{4}z_{153} + z_{6} - z_{152} - z_{153} + z_{154}\\*
&\quad{} + z_{155}z_{156}\\[4pt]
Q_{111} &={}1 - 2z_{5} + z_{5}^{2} - 2z_{5}z_{148} + z_{157}\\[4pt]
Q_{112} &={}-z_{5}z_{203} + z_{19} + z_{157}z_{158} - z_{202} + z_{204}\\[4pt]
Q_{113} &={}-1 + 2z_{148}z_{218} - z_{159} + 2z_{218} - z_{218}^{2}\\[4pt]
Q_{114} &={}-1 + z_{107} + z_{159}z_{160} - z_{202} - z_{203}z_{218} + z_{204}\\[4pt]
Q_{115} &={}-z_{148}z_{162} - z_{162} + z_{163}\\[4pt]
Q_{116} &={}2z_{161} + z_{161}^{2} + z_{162}^{2} - z_{163}^{2}\\[4pt]
Q_{117} &={}1 - 2z_{107} + z_{107}^{2} - 2z_{107}z_{148} + z_{164}\\[4pt]
Q_{118} &={}-1 - z_{107}z_{162} + z_{108} - z_{161} + z_{163} + z_{164}z_{165}\\[4pt]
Q_{119} &={}4 - 4z_{107} + z_{107}^{2} - 2z_{107}z_{148} + 2z_{148} + z_{166}\\[4pt]
Q_{120} &={}-1 - z_{107}z_{153} + z_{109} - z_{152} + z_{153} + z_{154}\\*
&\quad{} + z_{166}z_{167}\\[4pt]
Q_{121} &={}1 + z_{148} - z_{148}z_{169} - z_{169} + z_{170}\\[4pt]
Q_{122} &={}1 + 2z_{168} + z_{168}^{2} - 2z_{169} + z_{169}^{2} - z_{170}^{2}\\[4pt]
Q_{123} &={}-1 + z_{107} - z_{107}z_{169} + z_{113} + z_{164}z_{171} - z_{168}\\*
&\quad{} + z_{170}\\[4pt]
Q_{124} &={}-2z_{148} - z_{148}^{2} + z_{172}\\[4pt]
Q_{125} &={}1 - z_{107} - z_{107}z_{152} - z_{107}z_{154} + z_{114} + z_{148}z_{154}\\*
&\quad{} + z_{152} - z_{153}z_{172} + 2z_{154} + z_{166}z_{173}\\[4pt]
Q_{126} &={}-z_{148}z_{175} - z_{175} + z_{176}\\[4pt]
Q_{127} &={}2z_{174} + z_{174}^{2} + z_{175}^{2} - z_{176}^{2}\\[4pt]
Q_{128} &={}-1 - z_{107}z_{175} + z_{117} + z_{164}z_{177} - z_{174} + z_{176}\\[4pt]
Q_{129} &={}-z_{148}z_{179} - z_{179} + z_{180}\\[4pt]
Q_{130} &={}2z_{178} + z_{178}^{2} + z_{179}^{2} - z_{180}^{2}\\[4pt]
Q_{131} &={}-1 - z_{107}z_{179} + z_{120} + z_{164}z_{181} - z_{178} + z_{180}\\[4pt]
Q_{132} &={}-z_{148}z_{183} - z_{183} + z_{184}\\[4pt]
Q_{133} &={}2z_{182} + z_{182}^{2} + z_{183}^{2} - z_{184}^{2}\\[4pt]
Q_{134} &={}-1 - z_{107}z_{183} + z_{123} + z_{164}z_{185} - z_{182} + z_{184}\\[4pt]
Q_{135} &={}-z_{148}z_{187} - z_{187} + z_{188}\\[4pt]
Q_{136} &={}2z_{186} + z_{186}^{2} + z_{187}^{2} - z_{188}^{2}\\[4pt]
Q_{137} &={}-1 - z_{107}z_{187} + z_{126} + z_{164}z_{189} - z_{186} + z_{188}\\[4pt]
Q_{138} &={}-z_{148}z_{191} - z_{191} + z_{192}\\[4pt]
Q_{139} &={}2z_{190} + z_{190}^{2} + z_{191}^{2} - z_{192}^{2}\\[4pt]
Q_{140} &={}-1 - z_{107}z_{191} + z_{127} + z_{166}z_{193} - z_{190} + z_{191}\\*
&\quad{} + z_{192}\\[4pt]
Q_{141} &={}-z_{148}z_{195} - z_{195} + z_{196}\\[4pt]
Q_{142} &={}2z_{194} + z_{194}^{2} + z_{195}^{2} - z_{196}^{2}\\[4pt]
Q_{143} &={}-1 - z_{107}z_{195} + z_{130} + z_{164}z_{197} - z_{194} + z_{196}\\[4pt]
Q_{144} &={}-z_{148}z_{199} - z_{199} + z_{200}\\[4pt]
Q_{145} &={}2z_{198} + z_{198}^{2} + z_{199}^{2} - z_{200}^{2}\\[4pt]
Q_{146} &={}-1 - z_{107}z_{199} + z_{133} + z_{164}z_{201} - z_{198} + z_{200}\\[4pt]
Q_{147} &={}1 - z_{148}z_{203} - z_{203} + z_{204}\\[4pt]
Q_{148} &={}-1 + 2z_{202} + z_{202}^{2} + z_{203}^{2} - 2z_{204} - z_{204}^{2}\\[4pt]
Q_{149} &={}-z_{107}z_{203} + z_{134} + z_{166}z_{205} - z_{202} + z_{203} + z_{204}\\[4pt]
Q_{150} &={}-z_{148}z_{207} - z_{207} + z_{208}\\[4pt]
Q_{151} &={}2z_{206} + z_{206}^{2} + z_{207}^{2} - z_{208}^{2}\\[4pt]
Q_{152} &={}-1 - z_{107}z_{207} + z_{138} + z_{164}z_{209} - z_{206} + z_{208}\\[4pt]
Q_{153} &={}-z_{148}z_{211} - z_{211} + z_{212}\\[4pt]
Q_{154} &={}2z_{210} + z_{210}^{2} + z_{211}^{2} - z_{212}^{2}\\[4pt]
Q_{155} &={}-1 - z_{107}z_{211} + z_{141} + z_{164}z_{213} - z_{210} + z_{212}\\[4pt]
Q_{156} &={}-z_{148}z_{215} - z_{215} + z_{216}\\[4pt]
Q_{157} &={}2z_{214} + z_{214}^{2} + z_{215}^{2} - z_{216}^{2}\\[4pt]
Q_{158} &={}-1 - z_{107}z_{215} + z_{144} + z_{164}z_{217} - z_{214} + z_{216}\\[4pt]
Q_{159} &={}-6 - 2z_{0} - 3z_{0}^{2} - 2z_{0}z_{4} - z_{1}^{2} - z_{2}^{2} - 2z_{4}\\*
&\quad{} - 2z_{4}^{2} - z_{5}^{2} - z_{6}^{2} + 2z_{10} - z_{10}^{2} - z_{11}^{2} - z_{12}^{2}\\*
&\quad{} - z_{13}^{2} - z_{14}^{2} - z_{15}^{2} - z_{16}^{2} - z_{17}^{2} - z_{18}^{2} - z_{19}^{2}\\*
&\quad{} - z_{105}^{2} + 2z_{107} - 2z_{107}^{2} - z_{108}^{2} - z_{109}^{2} - z_{113}^{2} - z_{114}^{2}\\*
&\quad{} - z_{117}^{2} - z_{120}^{2} - z_{123}^{2} - z_{126}^{2} - z_{127}^{2} - z_{130}^{2} - z_{133}^{2}\\*
&\quad{} - z_{134}^{2} - z_{138}^{2} - z_{141}^{2} - z_{144}^{2} + z_{147} - z_{218}^{2}\\[4pt]
Q_{160} &={}-z_{5}z_{10} - z_{5}z_{11} - z_{5}z_{12} - z_{5}z_{13}\\*
&\quad{} - z_{5}z_{14} - z_{5}z_{15} - z_{5}z_{16} - z_{5}z_{17}\\*
&\quad{} - z_{5}z_{18} + z_{218}\\[4pt]
Q_{161} &={}-16 + z_{0} - z_{219}^{2}\\[4pt]
Q_{162} &={}-1 + z_{2} + z_{6} - z_{6}^{2}\\[4pt]
Q_{163} &={}-z_{148}z_{221} - z_{221} + z_{230}\\[4pt]
Q_{164} &={}-1 + z_{220}^{2} + z_{221}^{2} - z_{230}^{2}\\[4pt]
Q_{165} &={}-z_{223} - z_{223}z_{224} + z_{231}\\[4pt]
Q_{166} &={}-1 + z_{222}^{2} + z_{223}^{2} - z_{231}^{2}\\[4pt]
Q_{167} &={}-z_{215}^{2} + z_{232}\\[4pt]
Q_{168} &={}z_{221} - 2z_{226}z_{232} - z_{232}\\[4pt]
Q_{169} &={}-4z_{221} - 8z_{221}z_{225} + z_{224}\\[4pt]
Q_{170} &={}-z_{148} - z_{220}z_{227} + z_{224}\\[4pt]
Q_{171} &={}-1 - z_{214} - z_{220}z_{228} + z_{222}\\[4pt]
Q_{172} &={}-z_{105} - 4z_{215}z_{229} + z_{223}\\[4pt]
Q_{173} &={}2z_{0}^{2} + 2z_{0}z_{4} - 2z_{0}z_{153} - 2z_{0}z_{162} + z_{1}^{2}\\*
&\quad{} - 2z_{1}z_{203} + z_{2}^{2} - 2z_{2}z_{191} + z_{4}^{2} - 2z_{4}z_{162}\\*
&\quad{} + z_{10}^{2} - 2z_{10}z_{169} + z_{11}^{2} - 2z_{11}z_{175} + z_{12}^{2}\\*
&\quad{} - 2z_{12}z_{179} + z_{13}^{2} - 2z_{13}z_{183} + z_{14}^{2} - 2z_{14}z_{187}\\*
&\quad{} + z_{15}^{2} - 2z_{15}z_{195} + z_{16}^{2} - 2z_{16}z_{199} + z_{17}^{2}\\*
&\quad{} - 2z_{17}z_{207} + z_{18}^{2} - 2z_{18}z_{211} + z_{105}^{2} - 2z_{105}z_{215}\\*
&\quad{} - z_{148}z_{233} + z_{152}^{2} + z_{153}^{2} + z_{161}^{2} + z_{162}^{2} + z_{168}^{2}\\*
&\quad{} + z_{169}^{2} + z_{174}^{2} + z_{175}^{2} + z_{178}^{2} + z_{179}^{2} + z_{182}^{2} + z_{183}^{2}\\*
&\quad{} + z_{186}^{2} + z_{187}^{2} + z_{190}^{2} + z_{191}^{2} + z_{194}^{2} + z_{195}^{2} + z_{198}^{2}\\*
&\quad{} + z_{199}^{2} + z_{202}^{2} + z_{203}^{2} + z_{206}^{2} + z_{207}^{2} + z_{210}^{2} + z_{211}^{2}\\*
&\quad{} + z_{214}^{2} + z_{215}^{2}\\[4pt]
Q_{174} &={}-z_{109}z_{234} + z_{134}\\[4pt]
Q_{175} &={}-z_{114}z_{235} + z_{134}\\[4pt]
Q_{176} &={}-z_{127}z_{236} + z_{134}\\[4pt]
Q_{177} &={}-1 - z_{3}z_{5} + z_{5} + z_{6} - z_{237}^{2} - z_{238}^{2} - z_{239}\\*
&\quad{} - z_{239}^{2}\\[4pt]
Q_{178} &={}z_{3}z_{5} - z_{6} - z_{240}^{2} - z_{241}^{2} - z_{242} - z_{242}^{2}\\[4pt]
Q_{179} &={}-2 - z_{5} + z_{107} - z_{243}^{2} - z_{244}^{2} - z_{245} - z_{245}^{2}\\[4pt]
Q_{180} &={}-2 - z_{105}^{2} - z_{153}^{2} - z_{162}^{2} + 2z_{169} - z_{169}^{2} - z_{175}^{2}\\*
&\quad{} - z_{179}^{2} - z_{183}^{2} - z_{187}^{2} - z_{191}^{2} - z_{195}^{2} - z_{199}^{2} - z_{203}^{2}\\*
&\quad{} - z_{207}^{2} - z_{211}^{2} + z_{215}^{2} - z_{246}^{2} - z_{247}^{2} - z_{248} - z_{248}^{2}\\[4pt]
Q_{181} &={}-1 - 4z_{111}^{2} - 4z_{116}^{2} - 4z_{119}^{2} - 4z_{122}^{2} - 4z_{125}^{2} - 4z_{129}^{2}\\*
&\quad{} - 4z_{132}^{2} + z_{134}^{2} - 4z_{136}^{2} - 4z_{140}^{2} - 4z_{143}^{2} - 4z_{146}^{2} - z_{249}^{2}\\*
&\quad{} - z_{250}^{2} - z_{251} - z_{251}^{2}
\end{align*}
\endgroup


\end{document}
